\section{Experimental Setup}
\label{sec:setup}
\subsection{Sample and Model}
We use the human-verified portion of VulnGym v0.1.4 at dataset commit \texttt{cd69f7e}~\cite{vulngymdata}. Excluding pilot entries and their advisories leaves 338 entries from 156 advisories. The frozen sampling procedure retains one representative per advisory and selects 80 entries plus ten reserves using seed 20260929. Reserves are not evaluated. The sample contains 60 business-logic and 20 traditional-vulnerability entries, 130 entry-specific unique EP/CO target files, and 74 canonical repository--commit pairs across 13 repositories.

The sample is concentrated: Openclaw contributes 45 entries, n8n 12, and Flowise eight; the other ten repositories contribute 15 entries. Distinct advisories do not imply independent repositories. The public description is capped at 1,500 characters with the frozen truncation marker. The pre-existing \texttt{PATH\_HINT} flag labels 29 descriptions matching a case-insensitive source-path or filename/extension pattern; the remaining 51 are \texttt{NO\_PATH\_HINT}. This flag inspects the description alone, omitting the supplied title. We retain the frozen labels and refer to them as description-only flags.

A post hoc input audit matches each EP/CO target basename as a case-insensitive substring of the archived title. At least one basename matches in 47/80 entries: 12/29 in the description-hint group and 35/51 in the other group (Table~\ref{tab:input-hints}). A match need not identify a unique repository path, and not every match is proved to originate from an appended suffix. Neither description stratum represents input free of filename hints.
% Generated by generate_input_output_diagnostics.py; see input_output_provenance.json
\begin{table*}[t]
\centering\small
\caption{Description-only hint flags and target-basename matches in supplied entry titles.}\label{tab:input-hints}
\begin{adjustbox}{max width=\textwidth}\begin{tabular}{lrr}
\toprule
Frozen description flag & Entries & Title basename matches \\
\midrule
PATH\_HINT & 29 & 12 \\
NO\_PATH\_HINT & 51 & 35 \\
\bottomrule
\end{tabular}\end{adjustbox}
\par\vspace{3pt}\begin{minipage}{\textwidth}\footnotesize Matching is case-insensitive substring matching of a target basename, not full-path matching or an estimate of hint effects. The two original flags are retained; neither defines a hint-free input group.\end{minipage}
\end{table*}


Both studies use DeepSeek V4 Pro with high reasoning effort through Pi. Pi 0.85.1, Python 3.12.3, WSL2, embedding artifacts, dependency fingerprints, configuration, command, and timestamps are recorded with the study artifacts. The client records the model setting, but does not make provider generation deterministic or verify an identical hidden backend throughout collection. Fixed schedule and analysis seeds therefore coexist with stochastic model outcomes.

\subsection{Two Separately Frozen Studies}
Each study schedules 80 entries $\times$ three configurations $\times$ three repeats, or 720 unique inference units. Every arm has 240 outcomes. Repeats are averaged within an entry before paired comparisons; they are not 240 independent statistical samples.

S0 records the original text-tool interface, renderer, parser, source filters, and capped retrieval indexes. Its Agent source predicate admits 578 additional nonblank files excluded by Pack's directory/size/binary rules. Pack and Agent permit all 130 target files; RAG's 40,000-chunk index cap excludes two, in \texttt{entry-00183} and \texttt{entry-00359}. Pack's delivered prefix contains no target-file source, despite target eligibility. This last observation is an acquisition result, not an eligibility exclusion.

S1's scientific implementation and analysis were frozen before its formal collection, after S0 outcomes and implementation issues were known. It uses the shared manifest and renderer of Section~\ref{sec:method}, full indexes totaling 2,033,973 chunks, a native Agent interface, and schedule seed 20261003. The largest snapshot contains 56,867 chunks. Table~\ref{tab:configuration-changes} records the bundled changes. They include retrieval tie/numerical details, scheduling, output validation, and timeout semantics as well as source eligibility.
% Generated by paper/generate_revision_artifacts.py; see paper/provenance.json
\begin{table*}[t]
\centering
\caption{Implementation dimensions differing between S0 and S1. The two studies use the same known sample, not a single-factor intervention.}\label{tab:configuration-changes}
\small
\begin{adjustbox}{max width=\textwidth}\begin{tabular}{p{.15\textwidth}p{.35\textwidth}p{.40\textwidth}}
\toprule
Dimension & S0 & S1 \\
\midrule
Source/index & Arm-specific filters; RAG cap of 40,000 chunks & Shared authorization; every nonblank eligible chunk indexed \\
Agent interface & Text tool JSON and explicit experiment loop & Native schemas and Pi loop; multiple calls per assistant message \\
Final parsing & JSON extracted with permissive candidate validation & Bounded envelope; at most five valid positive integer intervals \\
Presentation & Pack tree; chunk interval headers; grep matches & Eligible path listing; numbered lines; grep context \\
Retrieval ties & Stable original ordering & Explicit sparse rank, path and starting line \\
Scheduling & Entry-first schedule & Repeat-first shuffled entry/arm schedule \\
Waiting timeout & 600 seconds per explicit prompt & 600 seconds, reset after each native assistant completion \\
\bottomrule
\end{tabular}\end{adjustbox}
\par\vspace{3pt}\begin{minipage}{\textwidth}\footnotesize S0 files are bound to the original launch manifest; S1 files are bound to the archived implementation. These descriptions are implementation observations, not estimated effects of individual changes.\end{minipage}
\end{table*}


Output validity is not identical across studies. Seven accepted S0 Agent answers contain a \texttt{line\_end=null} candidate that would invalidate the entire answer under S1's schema. Applying the strict schema to saved S0 answers gives Agent FR@5 0.8188 and LR@5 0.6417, with the original full denominator. This offline sensitivity check is not a replacement for S0 scores or a simulation of generation under S1 prompts. We retain each study's frozen output policy and do not pool its observations with the other study.

\subsection{Localization Metrics}
For entry $i$, let $T_i$ contain the one or two unique EP/CO paths and $P_{ir,5}$ the paths in the first five predictions for repeat $r$. The unit file score is
\begin{equation}
 \mathrm{FR@5}_{ir}=\frac{|T_i\cap P_{ir,5}|}{|T_i|}.
\end{equation}
The reported arm mean averages the three repeats per entry and then the 80 entries. It is not a pooled proportion of the 130 target files. It measures coverage rather than precision.

The shared user instruction ends with \emph{Identify the most likely vulnerable code locations.} The system instruction requests up to five likely vulnerable files with line intervals. Neither explicitly requests separate EP and CO predictions; the returned list is undifferentiated. Accordingly, the role-based scores measure coverage of benchmark annotations, not compliance with an explicit two-role task.

LR@5 keeps a denominator of two annotations, EP and CO, even when they share a path. For annotated interval $[u,v]$, a prediction hits if its path matches and its schema-valid integer interval $[a,b]$ satisfies
\begin{equation}
 a-5\leq v\quad\text{and}\quad b+5\geq u.
\end{equation}
An annotation without a qualifying prediction contributes zero. This criterion permits interval overlap with five-line tolerance; it does not require exact lines or containment of the full annotation. For example, a same-file prediction [1,1000] hits an annotation [500,510] despite its coarse extent. The frozen S1 supplementary metric LR@5$_{\leq64}$ additionally discards predicted intervals with inclusive width $b-a+1>64$; its denominator remains two and its hit rule remains overlap-based. We report it descriptively, without adding a post hoc test family.

\subsection{Analysis and Retained Failures}
The primary contrast frozen before collection is Agent--Pack FR@5. The secondary family comprises Agent--RAG and RAG--Pack contrasts for FR@5 and LR@5. Two-sided Monte Carlo sign-flip tests use 100,000 draws and a plus-one correction. Percentile intervals use 10,000 entry-bootstrap draws; analysis seed is 20260928. The primary threshold is $p<0.05$. Four secondary tests use Benjamini--Hochberg adjustment~\cite{bh} at $q<0.10$. S0 adjusts probabilities rounded to four decimals; S1 adjusts unrounded probabilities. Original statistics are retained.

Both procedures assume independence between entry-level differences. We therefore report repository means, leave-one-repository-out aggregates, an equal-repository mean, and the frozen repository-resampling sensitivity interval. Displayed entry intervals are unadjusted for both clustering and multiplicity.

During manuscript review, we add a separate post hoc sensitivity analysis for all five frozen contrasts. It enumerates all $2^{13}=8{,}192$ sign assignments to repository sums of the repeat-averaged entry differences. This retains the entry-weighted contrast as its estimand but assumes independently sign-exchangeable repository contributions. An exploratory four-secondary-test BH adjustment is calculated separately from the frozen entry-level family. Equal-repository means instead give every repository equal weight, and leave-one-repository-out means expose concentration. These diagnostics neither replace the official tests nor provide a general cluster-robust estimator; no CR2 analysis was implemented. We describe the original description-only hint strata using every outcome. Additional post hoc diagnostics audit title basenames and jointly classify source delivery and output at the annotation level. Each arm retains 480 opportunities: two annotations in each of 240 outcomes, even when EP and CO share a path. Delivered annotation source requires an actual recorded span to overlap the annotation, with zero tolerance. Prediction hits retain the frozen five-line tolerance. Invalid output takes precedence; other opportunities are partitioned by hit/miss, target-file delivery, and annotated-source delivery (Table~\ref{tab:annotation-stages}). These categories use the union of saved delivered spans and final predictions. They measure co-occurrence, not a temporal sequence or verified evidence use. We cross-tabulate them by annotation role and retention under the frozen 64-line rule.

We also report EP and CO recall separately, predicted-interval width quantiles, and the union of predicted lines within each file. Width quantiles describe valid returned intervals; mean union extent retains all 240 outcomes, assigning zero to failures. A paired decomposition partitions the Agent--RAG score difference by whether Agent returned a valid output, keeping 240 as the denominator for both contributions. We additionally check each top-five prediction against the full pinned Git tree and the file's length. Path existence is distinct from eligibility; existing filtered files are counted separately. The original union extent is model-reported numeric range size. A clipped extent removes nonexistent paths, intersects each interval with the actual file, and then merges within each path, averaging all 240 outcomes. Neither extent measures analyst effort. These diagnostics leave frozen scores and tests unchanged and do not estimate repaired-system performance.

An empty or invalid final model output scores FR@5=LR@5=0 and remains in the denominator. Completed output failures are never retried or replaced. Explicitly authorized infrastructure recovery resumes incomplete units while retaining interrupted attempts, diagnostics, and billing status. S1's completed matrix preserves all 31 Agent output failures and all 720 scheduled keys. The complete frozen schedule, code archive, recovery ancestry, delivered blocks, and effective-outcome fee journal are bound to the analysis artifacts; historical unresolved charges remain separate.
